\section{Threats to validity}\label{sec:threats}

\paragraph{Construct validity}
Labels follow the datasets' convention: a function is vulnerable if it was modified by a vulnerability-fixing change, and benign otherwise. Benign functions may contain undiscovered flaws, and vulnerable ones may be mislabelled; label noise in such corpora is documented~\cite{croft2023data}. Our ``miss rate'' is therefore the miss rate with respect to these labels, and the true miss rate in deployment could be higher; the guarantee is only as meaningful as the calibration labels, which is why the monitoring step of Algorithm~\ref{alg:triage} matters. The publication year of a CVE is a proxy for the age of the code, and the CWE annotation is incomplete, which limits the per-CWE analysis to classes with at least 15 vulnerable test functions.
We remove exact duplicates after normalisation but not near-duplicates (clones with renamed identifiers or minor edits), so random-split numbers are probably still optimistic; this makes the contrast with the shifted scenarios conservative.
The semantics-preserving rewrites are implemented with lexical analysis rather than a compiler front-end; renaming is restricted to declared locals and parameters and was not validated by recompilation.

\paragraph{Internal validity}
Hyper-parameters are fixed across scenarios instead of being tuned per scenario, so absolute detection scores are lower than what careful tuning might reach; the triage conclusions concern the \emph{relationship} between nominal and realised miss rates, which does not depend on the absolute score. The realised miss rates are estimated on roughly 700--11\,200 vulnerable test functions per scenario and seed in the main study (about 750 in the CodeBERT subsample), giving a sampling error of about one percentage point at $\alpha=10\%$ in the smallest case; we use a two-point tolerance for violation counts and report standard deviations over three seeds. The conformal implementation is checked by a unit test against the $p$-value definition and by simulation of the finite-sample guarantee. The domain classifier behind the weighted variant is trained on the same features as the detector, and a better density-ratio estimator could do better, so the mixed result for weighting should be read as ``not reliably effective with standard estimators'' rather than ``never effective''. The online experiment replays static test sets as streams with complete, delayed labels, uses a single scorer (LR+LGBM), and the step size was compared post hoc on the same streams, so it shows what the method can achieve when misses are discovered, not how fast they are discovered in practice.

\paragraph{External validity}
We study three C/C++ corpora that share provenance (43\% of Big-Vul and 77\% of PrimeVul also occur in DiverseVul, Tables~\ref{tab:data} and~\ref{tab:data3}), so the cross-dataset scenarios are harder than a random split but not independent tests. Because the main experiments were designed to run on CPUs, we evaluate four lightweight detectors on the full corpora, and CodeBERT, frozen and fine-tuned, only on a 30\,000-function stratified subsample with the first 256 tokens (which truncates $34$--$44\%$ of the normalised functions) and, for fine-tuning, on one or two seeds because the GPU quota ended; larger or newer code models and large language models were not evaluated and may reach higher AUROC, although fine-tuned CodeBERT did not outperform the lexical scorers here. We expect the logic of our findings to hold, as the CodeBERT experiment suggests, but the conformal guarantee depends on exchangeability and not on the scorer, and the \emph{possibility} of a violation under shift follows from Proposition~\ref{prop:shift} (a diagnostic upper bound), while its size must be measured empirically, as the range from near-nominal to $3.7\times$ the budget in our own scenarios shows. The \emph{magnitude} of the violation and the size of the cost-saving may differ for stronger models and should be re-measured; this is the most important item of future work. Results on other languages, on repository-level detection and on industrial data may differ.

The formatting shortcut we report was found in the Hugging Face releases of Big-Vul (clearly) and PrimeVul (more weakly, as function size alone is nearly as informative there) used here; we did not trace whether it originates in the original data collection or in the release, and because many PrimeVul benign functions are patched twins of vulnerable ones, the differences may partly reflect how the fixes were applied. Other releases may differ. The formatting-only probe is an upper-bound check, not an explanation of every model's behaviour.

\paragraph{Conclusion validity}
Blocks (scenario--seed pairs) are not fully independent because scenarios share data, so $p$-values are indicative; we therefore also report effect sizes and the full tables. The cost analysis depends on the assumed cost ratio $c_m$, which we vary over two orders of magnitude instead of picking one.
